\SourceLeaf{315}%
Once the value of \(T\) no longer changes when five-digit logarithms are used, \(\lambda\) is calculated from the formula
\[
\lambda=r\sin T\sec B
\]
A new calculation is then carried out to seven decimal places, using the logarithmic corrections from the formulas
\[
\begin{aligned}
\log b&=\log\beta+\mu\lambda\lambda+\tfrac12\mu\tau\tau\\
\log t&=\log\tau+\mu rr+\tfrac12\mu\tau\tau
\end{aligned}
\]
and setting \(B=B^0-\frac12b\) and \(T=T^0-\frac12t\). One further repetition will generally give the same results or results that differ only imperceptibly. Only then is the calculation of \(l\) added, using the following formula:
\[
\log l=\log\lambda-\tfrac12\mu rr+\tfrac12\mu\lambda\lambda
\]
To show by an example how quickly the approximation proceeds (which depends chiefly on how small \(r\) is), I give here the main steps of the calculation from the triangulation point Brocken to the point Inselsberg. This is the longest triangle side in the Hanoverian triangulation network, much longer than the sides usually encountered in trigonometric surveys.

In the conformal representation on the sphere worked out on the principles of the first memoir, the latitude of Brocken is \(B^0=51^0\,46'\,3''\,6345\), and the azimuth of the side Brocken-Inselsberg is \(T^0=5^0\,42'\,21''\,7704\). The logarithm of the length of this side is \(=4{,}7353929\) in toises, \(=8{,}22018543\) in units of the radius, or \(\log r=3{,}5346106\) in arcseconds, as assumed in our formulas. Setting \(T=5^0\,42'\) at first gives
\[
\begin{aligned}
\beta&=3408''\\
B&=51^0\,17'\,40''\\
\tau&=424''
\end{aligned}
\]
and hence a closer approximation
\[
T=5^0\,38'\,50''
\]
